% Partie : théorie de la démonstration
% Chapitre : déduction naturelle
% Section : séquents

\documentclass[../../../include/open-logic-section]{subfiles}

\begin{document}

\olfileid[fr]{pt}{nat}{seq}

\olsection{Déduction naturelle avec séquents : \Log{N2c} et~\Log{N2i}}

!!^a{proof}~$\delta$ dans \Log{N1c} ou \Log{N1i} établit
que sa !!{formula} finale~$!A$ découle de ses hypothèses
ouvertes. Si $\Gamma$ est l'ensemble des !!{formula}s~$!B$
telles que $!B^x$ apparaisse comme hypothèse ouverte
dans~$\delta$, on peut écrire $\delta \Proves \Gamma
\Sequent !A$. Cela suggère de présenter la déduction
naturelle au moyen de séquents. Chaque séquent indique
alors les hypothèses ouvertes dont dépend la formule
à droite : ce sont les hypothèses ouvertes de la
sous-!!{proof} qui se termine par ce séquent. Cette
présentation, appelée déduction naturelle avec séquents,
donne les systèmes \Log{N2c} et~\Log{N2i}.

Pour rapprocher les présentations \Log{N1} et \Log{N2},
prenons pour ensembles~$\Gamma$ à gauche des ensembles de
!!{formula}s étiquetées~$x:!B$. Les règles de \Log{N1}
n'agissent que sur les !!{formula}s prémisses, c'est-à-dire
les !!{formula}s finales des sous-!!{proof}s immédiates.
De même, les formules principales et actives de \Log{N2}
se trouvent dans le conséquent. À gauche, le \emph{contexte}
réunit les hypothèses étiquetées encore ouvertes.

Les étiquettes doivent être compatibles dans l'arbre entier :
une même étiquette ne désigne pas deux formules différentes.
Dans un schéma de décharge, le contexte résiduel ne
contient pas la classe étiquetée que l'on retire.

Dans les !!{proof}s de \Log{N2}, les séquents initiaux
remplacent les hypothèses ; ce sont des axiomes de la forme
\[x:!A \Sequent !A.\]
Les règles correspondent exactement à celles de \Log{N1}
et figurent dans le \olref{tab:N2}.

\subfile{rules-N2}

Comme dans les systèmes \Log{N1}, les règles
\Intro{\lif}, \Intro{\lnot}, \Elim{\lor}, \Elim{\lexists}
et, dans le système classique, \FalseCl peuvent s'appliquer
avec une décharge vacante. Dans les systèmes \Log{N2},
la classe à décharger peut donc être absente du contexte
de la prémisse concernée.



\begin{prop}
Si $\delta$ est une !!{proof} de~$!A$ dans~\Log{N1c}
(respectivement \Log{N1i}), il existe une
!!{proof}~$\delta'$ de $\Gamma \Sequent !A$ dans~\Log{N2c}
(respectivement \Log{N2i}), où $\Gamma$ est l'ensemble
des~$x:!B$ tels que $!B^x$ soit une hypothèse ouverte
de~$\delta$.
\end{prop}

\begin{proof}
Procédons par récurrence sur $n=\pheight{\delta}$.
Si $n=0$, $\delta$ est une unique hypothèse ouverte~$!A^x$.
Alors $\Gamma = \{x:!A\}$, et $\Gamma \Sequent !A$ est
un axiome de \Log{N2c} ou~\Log{N2i}.

Si $n>0$, distinguons les cas selon la dernière inférence
de~$\delta$. Nous en traitons un ; les autres sont laissés
en exercice.

Supposons que la dernière inférence soit $\Intro{\lif}$.
Alors $\delta$ se termine par
\[
  \AxiomC{$\Discharge{!B}{x}$}
  \RightLabel{$\delta_1$}
  \DeduceC{$!C$}
  \DischargeRule{\Intro{\lif}}{x}
  \UnaryInfC{$!B \lif !C$}
  \DisplayProof
  \]
L'hypothèse de récurrence donne une
!!{proof}~$\delta_1'$ de $\Gamma' \Sequent !C$ dans~\Log{N2c}
(respectivement \Log{N2i}), où $\Gamma'$ est l'ensemble des
hypothèses étiquetées ouvertes dans~$\delta_1$.
Si $!B^x$ est une hypothèse ouverte de $\delta_1$, elle
est déchargée par \Intro{\lif}. L'ensemble des hypothèses
ouvertes de~$\delta$ est donc
$\Gamma = \Gamma' \setminus \{x:!B\}$.
Ainsi, $\delta_1'$ est une !!{proof}
de~$x:!B, \Gamma \Sequent !C$, et l'on prend pour $\delta'$
\[
  \AxiomC{}
  \RightLabel{$\delta_1'$}
  \Deduce$x:!B, \Gamma \fCenter !C$
  \RightLabel{\Intro{\lif}}
  \UnaryInf$\Gamma \fCenter !B \lif !C$
  \DisplayProof
  \]
Si $!B^x$ n'est pas une hypothèse ouverte de~$\delta_1$,
l'inférence \Intro{\lif} ne décharge aucune hypothèse.
Les hypothèses ouvertes de $\delta$ et de $\delta_1$
sont alors les mêmes. Comme la formule de contexte
$x:!B$ n'est pas requise dans la prémisse de \Intro{\lif}
dans~\Log{N2c} ou~\Log{N2i}, on prend pour $\delta'$
\[
  \AxiomC{}
  \RightLabel{$\delta_1'$}
  \Deduce$\Gamma \fCenter !C$
  \RightLabel{\Intro{\lif}}
  \UnaryInf$\Gamma \fCenter !B \lif !C$
  \DisplayProof
  \]
Les autres dernières règles se traitent de manière analogue.
\end{proof}

\begin{prop}
Si $\delta$ est une !!{proof} de $\Gamma \Sequent !A$
dans \Log{N2c} (respectivement \Log{N2i}), il existe une
!!{proof}~$\delta'$ dans \Log{N1c} (respectivement \Log{N1i})
de !!{formula} finale~$!A$, dont les hypothèses ouvertes
sont exactement les $!B^x$ correspondant aux
$x:!B \in \Gamma$.
\end{prop}

\begin{proof}
Procédons encore par récurrence sur la !!{height}~$n$
de~$\delta$. Si $n=0$, $\delta$ est un axiome
$x:!A \Sequent !A$, et la !!{proof}~$\delta'$ est
l'hypothèse~$!A^x$.

Si $n>0$, distinguons les cas selon la dernière inférence.
Par exemple, si c'est~$\Intro{\lif}$, $\delta$ se termine par
\[
  \bottomAlignProof
  \AxiomC{}
  \RightLabel{$\delta_1$}
  \Deduce$x:!B, \Gamma \fCenter !C$
  \RightLabel{\Intro{\lif}}
  \UnaryInf$\Gamma \fCenter !B \lif !C$
  \DisplayProof
  \quad\text{ ou }\quad
  \bottomAlignProof
  \AxiomC{}
  \RightLabel{$\delta_1$}
  \Deduce$\Gamma \fCenter !C$
  \RightLabel{\Intro{\lif}}
  \UnaryInf$\Gamma \fCenter !B \lif !C$
  \DisplayProof
  \]
L'hypothèse de récurrence donne une
!!{proof}~$\delta_1'$ de~$!C$ dans~\Log{N1c}
(respectivement \Log{N1i}), avec exactement les hypothèses
ouvertes étiquetées du contexte de la prémisse.
Dans le premier cas, $!B^x$ est une hypothèse ouverte de
cette sous-preuve ; dans le second, elle ne l'est pas.
On obtient la !!{proof}~$\delta'$ en appliquant
\Intro{\lif}, avec décharge dans le premier cas et
sans décharge dans le second :
\[
  \bottomAlignProof
  \AxiomC{$\Discharge{!B}{x}$}
  \RightLabel{$\delta_1'$}
  \DeduceC{$!C$}
  \DischargeRule{\Intro{\lif}}{x}
  \UnaryInfC{$!B \lif !C$}
  \DisplayProof
  \quad\text{ ou }\quad
  \bottomAlignProof
  \AxiomC{}
  \RightLabel{$\delta_1'$}
  \DeduceC{$!C$}
  \RightLabel{\Intro{\lif}}
  \UnaryInfC{$!B \lif !C$}
  \DisplayProof
  \]
Les autres règles se traitent de même ; pour les règles
de quantification, le contexte enregistre précisément
les hypothèses ouvertes auxquelles s'applique la
condition de la constante propre.
\end{proof}












\end{document}
